\chapter{Introduction}

\section{Motivation}
Courier companies must know where each parcel is at any time. This project
builds a small simulator of that task. It lets the group practice the main
ideas of the Object-Oriented Programming Lab in one working program.

\section{Background}
A parcel moves through four hubs in a fixed order: Dispatch, Sorting Facility,
Out-for-Delivery, and Delivered. Each move is a milestone (one event in the
history of a parcel). The program saves every milestone in a text file.

The program uses these object-oriented ideas: classes and objects,
inheritance, abstract methods, polymorphism, exception handling, and file
input and output.

\section{Summary of present work}
The group built a command-line program in core Java. The user can add a
parcel, move it to the next hub, and see its tracking history. There are three
shipment speeds: Standard, Express, and Overnight. The program rejects invalid
tracking IDs with a custom exception. It saves the parcels and the events in
text files, so the data is still there after the program restarts.

\section{Organization of the report}
Chapter 2 states the problem and the requirements. Chapter 3 describes the
design. Chapter 4 explains the implementation. Chapter 5 presents testing and
results. Chapter 6 gives the conclusions and Chapter 7 gives the future scope.

\section{Hardware/Software used}
The program uses core Java only. It uses no database and no outside library.
The group wrote the code in Visual Studio Code and kept the project in a
GitHub repository. The tests are Java classes that one script runs. The Java documentation
and tutorials \cite{ref1,ref2,ref3} were the main references.
